\scanpage{22}

% L1–4
du produit 2-fibré dans (Top) $[\,\deleted{\unsure{de|des}}\,R\mathop{\times}^{2}_{S}R_1\,]$, ou la 2-somme amalgamée dans $\Delta$ \quad $R\amalg_S R_1$. Cette construction est $\pm$ connue, redonnons en une autre, dans le cas général $R\mathop{\times}^{2}_{T}S$, avec les foncteurs images inverses
% L4, right
\diagram{\begin{tikzcd}[ampersand replacement=\&,column sep=small,row sep=tiny]
 {} \& R\\ T\arrow[ur,"f^*"]\arrow[dr,"g^*"'] \& {}\\ {} \& S
\end{tikzcd}.}
% L4–7
On choisit des \deleted{familles} \inserted{petites sous-catégories} génératrices $T_0,R_0,S_0$ telles que $f^*(T_0)\subset R_0$, $g^{\unsure{*|'}}(T_0)\subset S_0$. On prend $U_0$ la catégorie-somme amalgamée de $R$ et $S$, de sorte que pour tte \inserted{catégorie} $E$ \deleted{\illegible{1}} \ldots .
% L8
\[
 \underline{\mathrm{Hom}}(U_0,E)\xrightarrow{\sim}
 \underline{\mathrm{Hom}}_{\deleted{\unsure{\Delta}}}(R_0,E)
 \mathop{\times}^{2}_{\underline{\mathrm{Hom}}(T_0,E)}
 \underline{\mathrm{Hom}}(S_0,E)
\]

% L9–14
Il faut, si $E\in\operatorname{Ob}\Delta$, exprimer que \unsure{l'un|un} objet du \deleted{second \unsure{membre|mb.}} premier mb, $U_0\xrightarrow{h}E$, est dans \deleted{$\underline{\mathrm{Hom}}_\Delta(R,E)$} la sous-catégorie pleine $\underline{\mathrm{Hom}}_\Delta(R,E)\mathop{\times}^{2}_{\underline{\mathrm{Hom}}_\Delta(T,E)}\underline{\mathrm{Hom}}_\Delta(\unsure{T|S},E)$ du deuxième mb. Si on a pris $R_0,S_0$ stables par \uline{lim} finies, cela s'exprime aisément par la \deleted{commutation de $h$ aux} condition que $h$ transforme certains cônes projectifs finis de $U_0$ (\inserted{\unsure{ceux|aux}} provenant des cônes exacts \deleted{de} $R_0$ et de $S_0$) en cônes exacts,
% L15–19
et \deleted{que les \unsure{im|uns}} qu'elle transforme certaines familles (provenant de familles couvrantes de $R_0$ et $S_0$) en familles couvrantes. Donc on réalise la théorie envisagée $\tau=T_R\mathop{\times}^{(2)}_{T_T}T_S$ comme sous-théorie pleine (définie par des axiomes d'\deleted{\unsure{égalité}}isomorphie, finitaires) de $T_{\rho_\Delta(U_0)}$, ce qui donne la conclusion.

% L19–26
Mais sachant désormais que \deleted{\unsure{Dans|dans}} la théorie $\tau$ est \unsure{représent.ble|représentable} par un topos, \deleted{et} revenant à la description générale comme sous-théorie d'un $T_{\rho_\Delta(U_0)}$ (sans \inserted{plus de} condition que $R_0$, $S_0$ stables par $\underleftarrow{\lim}$ existantes, on voit que c'est donné par un \uline{sous}-\uline{topos} de $\rho_\Delta(U_0)$. Ainsi on trouve, si $T_0\to S_0$ est surjectif sur les objets (comme dans le cas qui nous intéressait) que $R_0-U_0$ est surjectif sur les objets, donc $R\to U=R'$ \deleted{est géné} a une image qui est génératrice.
% archive sheet number: 21
